% Awesome Source CV LaTeX Template
%
% This template has been downloaded from:
% https://github.com/darwiin/awesome-neue-latex-cv
%
% Author:
% Christophe Roger
%
% Template license:
% CC BY-SA 4.0 (https://creativecommons.org/licenses/by-sa/4.0/)
%Section: Project
\IfLanguageName{english}{
    \sectionTitle{Projects}{\faLaptop}

    \begin{projects}

        \ifcvshowprojrenderer
        \project
    	{Renderer}{2020-now}
    	{\github{GraphicsEngineVulkan} \github{GraphicEngine}}
        {This projects provide me with a solid Vulkan/OpenGL starting point for implementing modern established rendering techniques and getting quickly started in own research topics.}
    	{CMake, C, C++, Python, Vulkan, OpenGL}
        \fi

        \ifcvshowprojhub
        \project
    	{Build \& CI Infrastructure}{2025-now}
    	{\github{ANTfrastructure}}
    	% Two lines since genesis-sim needed the room; the three-line version, kept: The build, test and CI hub my other public repositories share: staged container images for Linux (x86-64, arm64, riscv64) and Windows (x64, arm64 cross), CUDA and ROCm variants, compiler caching end to end, digest-pinned stage handoffs, and mutation-tested quality gates sharded across eight CI jobs.
    	{The build, test and CI hub my public repositories share: multi-arch container images for Linux and Windows with CUDA and ROCm variants, compiler caching and mutation-tested quality gates sharded across eight CI jobs.}
    	{Python, Bash, PowerShell, BuildKit, sccache, GitHub Actions}
        \fi

        \ifcvshowprojwgpu
        \project
    	{WebGPU Renderer in Rust}{2026-now}
    	{\github{OxidANT}}
    	% Metal dropped until it is tested on a Mac (2026-11-15); the claim before, kept: runs natively (Vulkan, DX12, Metal)
    	{wgpu glTF renderer that runs natively (Vulkan, DX12) and in the browser (wasm32, WebGPU): PBR with image-based lighting, cascaded shadows, SSAO, bloom, GPU skinning, an egui overlay and headless golden-image tests. It shares its Slang shaders with my C++ Vulkan engine.}
    	{Rust, wgpu, WGSL, Slang, egui, WebAssembly}
        \fi

        \ifcvshowprojomni
        \project
    	{Cross-Platform Inference App}{2024-now}
    	{\github{OmniAccelerANT}}
    	{Flutter frontend on a Rust and C++ core through flutter\_rust\_bridge: GStreamer camera pipelines, WebRTC live streams to the browser and ONNX inference. Its cat cam, packaged for Linux and Windows, logs every detected visit with a picture and a 10\,s clip.}
    	{Flutter, Dart, Rust, C++, GStreamer, WebRTC}
        \fi

        \ifcvshowprojml
        \project
    	{Machine Learning Algorithms}{2020 - now}
    	{\github{MachineLearningAlgorithms}}
    	{This project provides me a solid Machine Learning(ML) starting point for implementing modern established ML algorithms and getting quickly started in own research topics.}
    	{\LaTeX, Python, R}
        \fi

        \ifcvshowprojmr
        \project
    	{User-Adaptive MR Guidance}{2022-2023}
    	{\github{Designing-User-adaptive-Content-for-Mixed-Reality-Using-Eye-and-Hand-Tracking}}
    	{M.Sc.\ thesis: Investigates interrelationships between users' fixated objects, eye/head/hand movement, and intentions. Provides user-adaptive guidance for improved interaction quality in Mixed Reality.}
    	{\LaTeX, C\sharp, Mixed Reality, Python, Unity}
        \fi

        \ifcvshowprojbluenoise
        \project
    	{Blue Noise Error Distribution}{2019-2020}
    	{\github{BachelorArbeit}}
    	{B.Sc.\ thesis: Temporally stable blue noise error distribution in screen space for real-time rendering applications, building on prior work on blue noise error distributions for image quality.}
    	{\LaTeX, C, C++, Python}
        \fi

        \ifcvshowprojcollision
        \project
    	{GPU Collision Detection}{2021-2022}
    	{\github{GPUCollisionDetectionWithAS}}
    	{Novel collision detection using GPU acceleration structures. Implementation of a particle system accelerated on the GPU with BVH-based broad phase.}
    	{\LaTeX, C, C++, Vulkan API, GPGPU}
        \fi

    \end{projects}
}{
    \sectionTitle{Projekte}{\faLaptop}

    \begin{projects}
        \ifcvshowprojrenderer
        \project
    	{Renderer}{2020-now}
    	{\github{GraphicsEngineVulkan}\github{GraphicEngine}}
    	{Diese Projekte bieten mir einen soliden Startpunkt um diverse etablierte Computergrafik Algorithmen nachzuimplementieren und darauf aufbauend auch eigene Forschung zu betreiben.}
    	{\LaTeX, CMake, C, C++, Python, Vulkan, OpenGL}
        \fi

        \ifcvshowprojhub
        \project
     	{Build- \& CI-Infrastruktur}{2025-now}
     	{\github{ANTfrastructure}}
     	% Gekürzt wie der englische Zweig; die Fassung davor, erhalten: Der Build-, Test- und CI-Hub, den meine anderen öffentlichen Repositories teilen: gestufte Container-Images für Linux (x86-64, arm64, riscv64) und Windows (x64, arm64-Cross), CUDA- und ROCm-Varianten, durchgängiges Compiler-Caching, per Digest gepinnte Stufenübergaben und mutationsgetestete Qualitäts-Gates, verteilt auf acht CI-Jobs.
     	{Der Build-, Test- und CI-Hub, den meine öffentlichen Repositories teilen: Multi-Arch-Container-Images für Linux und Windows mit CUDA- und ROCm-Varianten, Compiler-Caching und mutationsgetestete Qualitäts-Gates, verteilt auf acht CI-Jobs.}
     	{Python, Bash, PowerShell, BuildKit, sccache, GitHub Actions}
        \fi

        \ifcvshowprojwgpu
        \project
     	{WebGPU-Renderer in Rust}{2026-now}
     	{\github{OxidANT}}
     	% Metal gestrichen, bis es auf einem Mac getestet ist (2026-11-15); vorher: nativ (Vulkan, DX12, Metal)
     	{wgpu-glTF-Renderer, der nativ (Vulkan, DX12) und im Browser (wasm32, WebGPU) läuft: PBR mit Image-Based Lighting, Cascaded Shadows, SSAO, Bloom, GPU-Skinning, egui-Overlay und Headless-Golden-Image-Tests. Seine Slang-Shader teilt er mit meiner C++-Vulkan-Engine.}
     	{Rust, wgpu, WGSL, Slang, egui, WebAssembly}
        \fi

        \ifcvshowprojomni
        \project
     	{Plattformübergreifende Inferenz-App}{2024-now}
     	{\github{OmniAccelerANT}}
     	{Flutter-Frontend auf einem Rust- und C++-Kern über flutter\_rust\_bridge: GStreamer-Kamera-Pipelines, WebRTC-Livestreams in den Browser und ONNX-Inferenz. Die für Linux und Windows paketierte Katzenkamera protokolliert jeden erkannten Besuch mit Bild und 10-s-Clip.}
     	{Flutter, Dart, Rust, C++, GStreamer, WebRTC}
        \fi

        \ifcvshowprojml
     	\project
     	{Machine Learning Algorithms}{2020 - now}
     	{\github{MachineLearningAlgorithms}}
     	{Dieses Projekt bietet mir einen soliden Startpunkt um diverse etablierte Machine Learning(ML) Algorithmen nachzuimplementieren und darauf aufbauend auch eigene Forschung zu betreiben.}
     	{\LaTeX, Python, R}
        \fi

        \ifcvshowprojmr
        \project
     	{Benutzeradaptive MR-Führung}{2022-2023}
     	{\github{Designing-User-adaptive-Content-for-Mixed-Reality-Using-Eye-and-Hand-Tracking}}
     	{Masterarbeit: Untersucht Wechselbeziehungen zwischen fixierten Objekten, Augen-/Kopf-/Handbewegungen und Absichten der Benutzer. Bietet benutzeradaptive Führung für verbesserte Interaktion in Mixed Reality.}
     	{\LaTeX, C\sharp, Mixed Reality, Python, Unity}
        \fi

        \ifcvshowprojbluenoise
        \project
     	{Blue Noise Fehlerverteilung}{2019-2020}
     	{\github{BachelorArbeit}}
     	{Bachelorarbeit: Zeitlich stabile Blue Noise Fehlerverteilung im Bildraum für Echtzeitanwendungen, aufbauend auf bisherigen Arbeiten zur Bildqualitätssteigerung.}
     	{\LaTeX, C, C++, Python}
        \fi

        \ifcvshowprojcollision
        \project
     	{GPU-Kollisionsdetektion}{2021-2022}
     	{\github{GPUCollisionDetectionWithAS}}
     	{Neuartige Kollisionsdetektion mit GPU-Beschleunigungsstrukturen. Implementierung eines GPU-beschleunigten Partikelsystems mit BVH-basierter Grobphase.}
     	{\LaTeX, C, C++, Vulkan API, GPGPU}
        \fi
        
    \end{projects}
}
